% ============ 4 平展拓扑的局部环 (原书 pp.31-37) ============
\chapter{平展拓扑的局部环}\label{sec:4}
\index{平展拓扑的局部环}
\index{Hensel 环}
\index{Hensel 化}
\index{严格 Hensel 化}
\index{严格局部环}
\index{几何点}
\index{平展邻域}
\index{Artin 逼近定理}

设 $X$ 是赋予函数层 $\mathcal{O}_X$ 的拓扑空间.函数在 $x$ 处的芽由一对 $(U, f)$ 表示，其中 $U$ 是 $x$ 的开邻域，$f$ 是 $U$ 上的函数；若对包含在 $U \cap U'$ 中的 $x$ 的某个开邻域 $V$，$f|_{V} = f'|_{V}$，则两对 $(U, f)$ 与 $(U', f')$ 表示同一个芽.换言之，函数在 $x$ 处的芽是环
\[
\mathcal{O}_{X, x} \overset{\mathrm{def}}{=} \varinjlim_{U} \Gamma(U, \mathcal{O}_X)
\]
的元素，其中 $U$ 遍历 $x$ 的开邻域.

例如，若 $X = \mathbb{C}$ 赋予全纯函数层，则函数在 $c \in \mathbb{C}$ 处的芽定义幂级数 $\sum_{n \geq 0} a_n (z - c)^{n}$，它在 $c$ 的某个开邻域中收敛，且每个这样的幂级数都来自唯一的芽；因此 $\mathcal{O}_{X, c}$ 是这种幂级数组成的环.若 $X$ 是代数闭域上的代数簇，则对包含 $x$ 的每个开仿射子簇 $U$，$\mathcal{O}_{X, x} = A_{\mathfrak{m}}$，其中 $A = \Gamma(U, \mathcal{O}_X)$，$\mathfrak{m}$ 是对应于 $x$ 的 $A$ 中的极大理想---见 AG 3.6.

更一般地，对任意环化空间 $(X, \mathcal{O}_X)$($\mathcal{O}_X$ 不一定是 $X$ 上的函数层)与 $x \in X$，我们定义 $\mathcal{O}_{X, x} = \varinjlim_{U} \Gamma(U, \mathcal{O}_X)$，其中 $U$ 遍历 $x$ 的开邻域.例如，若 $X$ 是概形，则对包含 $x$ 的任意开仿射 $U$，$\mathcal{O}_{X, x} = A_{\mathfrak{p}}$，其中 $A = \Gamma(U, \mathcal{O}_X)$，$\mathfrak{p}$ 是对应于 $x$ 的 $A$ 中的素理想.

局部环化空间是环化空间 $(X, \mathcal{O}_X)$，使得对所有 $x \in X$，$\mathcal{O}_{X, x}$ 都是局部环---这时 $\mathcal{O}_{X, x}$ 称为 $x$ 处(对所给环化空间结构)的局部环.我们将研究平展拓扑的类似环.

\subsection*{簇的情形}

本节中，$X$ 是代数闭域 $k$ 上的簇.

回忆点 $x \in X$ 的平展邻域是平展映射 $U \to X$ 连同映到 $x$ 的点 $u \in U$.平展邻域的态射(或映射) $(V, v) \to (U, u)$ 是把 $v$ 送到 $u$ 的 $X$ 上的正则映射 $V \to U$.由推论~\ref{cor:2.16} 推出，从连通的平展邻域到另一个平展邻域至多有一个映射.连通仿射平展邻域构成有向集\footnote{实际上，它们不构成集合---只是一个类.然而，可以用集合替换这个类.例如，若 $V$ 仿射(我们可以这样假定)，设 $V = \operatorname{Specm}(A)$，那么我们只需考虑平展映射 $\operatorname{Specm} B \to \operatorname{Specm} A$，其中 $B$ 是形如 $A[T_1, \ldots, T_n]$(某个固定的符号集 $T_1, T_2, \ldots$)的 $A$-代数的商的环的分式环.类似的评注也适用于别处，通常不再赘述.}，其定义为
\[
(U, u) \leq (U', u') \quad \text{若存在映射 } (U', u') \to (U, u),
\]
我们定义 $x$ 处平展拓扑的局部环为
\[
\mathcal{O}_{X, \bar{x}} = \varinjlim_{(U, u)} \Gamma(U, \mathcal{O}_U).
\]
由于 $x$ 的每个 Zariski 开邻域也是 $x$ 的平展邻域，我们得到同态 $\mathcal{O}_{X, x} \to \mathcal{O}_{X, \bar{x}}$.类似地，对 $x$ 的任意平展邻域 $(U, u)$，得到同态 $\mathcal{O}_{U, u} \to \mathcal{O}_{X, \bar{x}}$，并且显然
\[
\mathcal{O}_{X, \bar{x}} = \varinjlim_{(U, u)} \mathcal{O}_{U, u}.
\]
正向系统中的转移映射 $\mathcal{O}_{U, u} \to \mathcal{O}_{V, v}$ 都是平坦(从而是单射)非分歧的局部环局部同态，且 Krull 维数为 $\dim X$.

\begin{proposition}\label{prop:4.1}
环 $\mathcal{O}_{X, \bar{x}}$ 是 Krull 维数为 $\dim X$ 的局部 Noether 环.
\end{proposition}

\begin{proof}
局部环的局部同态系统的正向极限是局部的(极大理想是诸极大理想的极限)；因此 $\mathcal{O}_{X, \bar{x}}$ 是局部的.完备化上的映射 $\widehat{\mathcal{O}}_{U, u} \to \widehat{\mathcal{O}}_{V, v}$ 都是同构---见命题~\ref{prop:2.9} 的证明---由此推出 $\widehat{\mathcal{O}}_{X, \bar{x}} = \widehat{\mathcal{O}}_{X, x}$.Nagata 的一个论证(Artin 1962, 4.2)现在表明 $\mathcal{O}_{X, \bar{x}}$ 是 Noether 的，从而 Krull 维数为 $\dim X$.
\end{proof}

\textbf{环 $\mathcal{O}_{X, \bar{x}}$ 是 Hensel 的} $\mathcal{O}_{X, \bar{x}}$ 最引人注目的性质是它满足 Hensel 引理的结论---事实上(见下一小节)，它的特征恰为：它是包含 $\mathcal{O}_{X, x}$ 的具有此性质的``最小''局部环.

\begin{definition}\label{def:4.2}
局部环 $A$ 称为 Hensel 的，如果它满足如下条件：
设 $f(T)$ 是 $A[T]$ 中的首一多项式，$\bar{f}(T)$ 表示它在 $k[T]$ 中的像($k = A/\mathfrak{m}_{A}$)；则 $\bar{f}$ 分解为两个互素的首一多项式的乘积的每个分解都提升为 $f$ 分解为两个首一多项式的乘积，即若 $\bar{f} = g_0h_0$，其中 $g_0$ 与 $h_0$ 在 $k[T]$ 中首一且互素，则 $f = gh$，其中 $g$ 与 $h$ 首一，$\bar{g} = g_0$ 且 $\bar{h} = h_0$.
\end{definition}

\begin{remark}\label{rem:4.3}
(a) 上述分解中的 $g$ 与 $h$ 严格互素，即 $A[T] = (g, h)$.要看出这一点，注意 $M \overset{\mathrm{def}}{=} A[T]/(g, h)$ 作为 $A$-模有限生成(因为 $g(T)$ 首一)，且 $M = \mathfrak{m}_{A}M$(因为 $(\bar{g}, \bar{h}) = k[T]$).现在 Nakayama 引理蕴含 $M = 0$.

(b) 上述分解中的 $g$ 与 $h$ 是唯一的，因为设 $f = gh = g'h'$，其中 $g, h, g', h'$ 都首一且 $\bar{g} = \bar{g'}$, $\bar{h} = \bar{h'}$.(a) 中的论证表明 $g$ 与 $h'$ 严格互素，因此存在 $r, s \in A[T]$ 使得 $gr + h's = 1$.现在
\[
g' = g'gr + g'h's = g'gr + ghs,
\]
从而 $g$ 整除 $g'$.由于它们次数相同且首一，它们必相等.
\end{remark}

\begin{theorem}\label{thm:4.4}
每个完备局部环都是 Hensel 的.
\end{theorem}

\begin{proof}
数论中 Hensel 引理的标准证明对任意完备局部环都成立---见 Atiyah 和 MacDonald 1969, p.115.
\end{proof}

\begin{theorem}\label{thm:4.5}
对 $X$ 中任意点 $x$，$\mathcal{O}_{X, \bar{x}}$ 是 Hensel 的.
\end{theorem}

这由下面两个引理推出.

考虑局部环 $A$ 的如下条件(同样用 $f \mapsto \bar{f}$ 表示商映射 $A[T] \to k[T]$，$k = A/\mathfrak{m}_{A}$):

\begin{quote}
($\ast$) 设 $f_1, \ldots, f_n \in A[T_1, \ldots, T_n]$；$\bar{f_i}$ 在 $k^{n}$ 中的每个使得 $\operatorname{Jac}(f_1, \ldots, f_n)(a_0)$ 非零的公共零点 $a_0$ 都提升为 $f_i$ 在 $A^{n}$ 中的公共零点.
\end{quote}

\begin{lemma}\label{lem:4.6}
对 $X$ 中任意点 $x$，$\mathcal{O}_{X, \bar{x}}$ 满足 ($\ast$).
\end{lemma}

\begin{proof}
设 $f_1, f_2, \ldots, f_n \in \mathcal{O}_{X, \bar{x}}[T_1, \ldots, T_n]$，$a_0 \in k^{n}$ 如 ($\ast$) 中所述.按定义
\[
\mathcal{O}_{X, \bar{x}} = \varinjlim k[U], \quad \mathfrak{m}_{\bar{x}} = \varinjlim \mathfrak{m}_{u}
\]
($x$ 的平展邻域 $(U, u)$ 上的正向极限；$k[U] = \Gamma(U, \mathcal{O}_U)$ 是 $U$ 上正则函数的环；$\mathfrak{m}_{\bar{x}}$ 是 $\mathcal{O}_{X, \bar{x}}$ 的极大理想；$\mathfrak{m}_{u}$ 是 $u$ 处为零的 $U$ 上正则函数的理想).因为只有有限多个 $f_i$，且每个只有有限多个系数，存在 $(U, u)$ 使得 $f_i \in B[T_1, \ldots, T_n]$，$B = k[U]$.

设 $C = B[T_1, \ldots, T_n]/(f_1, \ldots, f_n)$，并设 $V = \operatorname{Specm}(C)$.同态 $B \to C$ 定义映射 $V \to U$.由 $a_0$ 得到位于 $u$ 之上的点 $v \in V$，而条件 $\operatorname{Jac}(f_1, \ldots, f_n)(a_0) \neq 0$ 蕴含 $V \to U$ 在 $v$ 处平展(见命题~\ref{prop:2.1}).因此，由命题~\ref{prop:2.14}，存在 $g \in C$，$g \notin \mathfrak{m}_{v}$，使得 $D(g) \to U$ 平展\footnote{记号如 AG：$D(g)$ 是 $g$ 非零的开集.}.于是存在同态 $C_g \to \mathcal{O}_{X, \bar{x}}$，使得 $\mathcal{O}_{X, \bar{x}}$ 的极大理想的原像是 $\mathfrak{m}_{v}$.但给出这样的同态就是给出 $f_i$ 在 $\mathcal{O}_{X, \bar{x}}$ 中提升 $a_0$ 的公共零点.
\end{proof}

\begin{lemma}\label{lem:4.7}
设 $A$ 是局部环.若 $A$ 满足 ($\ast$)，则它是 Hensel 的.
\end{lemma}

\begin{proof}
设
\[
f(T) = T^{n} + a_1 T^{n-1} + \cdots + a_n,
\]
并假定 $\bar{f}(T) = \bar{g_0}(T)\bar{h_0}(T)$，其中 $\bar{g_0}$ 与 $\bar{h_0}$ 次数分别为 $r$ 与 $s$，且互素.我们寻求分解
\[
f(T) = (T^{r} + x_1 T^{r-1} + \cdots + x_r)(T^{s} + y_1 T^{s-1} + \cdots + y_s)
\]
于 $A[T]$ 中，提升 $\bar{f}(T)$ 的给定分解.为得到分解，我们必须解方程组
\[
\begin{array}{c}
x_1 + y_1 = a_1 \\
x_2 + x_1 y_1 + y_2 = a_2 \\
x_3 + x_2 y_1 + x_1 y_2 + y_3 = a_3 \\
\vdots \\
x_r y_s = a_n
\end{array}
\]
($n$ 个方程，$n$ 个未知数).分解 $\bar{f} = \bar{g_0}\bar{h_0}$ 提供该方程组模 $\mathfrak{m}_{A}$ 的解.方程组的 Jacobi 矩阵是
\[
\begin{pmatrix}
1 & 0 & 0 & \cdots & 0 & 1 & 0 & 0 & \cdots & 0\\
y_1 & 1 & 0 & \cdots & 0 & x_1 & 1 & 0 & \cdots & 0\\
y_2 & y_1 & 1 & \cdots & 0 & x_2 & x_1 & 1 & \cdots & 0\\
\vdots & \vdots & \vdots & \ddots & \vdots & \vdots & \vdots & \vdots & \ddots & \vdots\\
y_s & y_{s-1} & y_{s-2} & \cdots & 1 & x_s & x_{s-1} & x_{s-2} & \cdots & 1\\
0 & y_s & y_{s-1} & \cdots & y_1 & 0 & x_s & x_{s-1} & \cdots & x_1\\
0 & 0 & y_s & \cdots & y_2 & 0 & 0 & x_s & \cdots & x_2\\
\vdots & \vdots & \vdots & \ddots & \vdots & \vdots & \vdots & \vdots & \ddots & \vdots\\
0 & 0 & 0 & \cdots & y_s & 0 & 0 & 0 & \cdots & x_s
\end{pmatrix}
\]
这个矩阵的转置的行列式是 $T^{r} + y_1 T^{r-1} + \cdots + y_r$ 与 $T^{s} + x_1 T^{s-1} + \cdots + x_s$ 的结式.因为 $\bar{g_0}$ 与 $\bar{h_0}$ 首一且互素，$\operatorname{Res}(\bar{g_0}, \bar{h_0}) \neq 0$(AG 7.11).因此，在分解 $\bar{f} = \bar{g_0}\bar{h_0}$ 提供的 $k$ 中的解处，方程组的 Jacobi 矩阵非零.现在 ($\ast$) 表明分解提升到 $A$.
\end{proof}

\textbf{非奇异点处的局部环} 微分流形上的点处的局部环只依赖于流形的维数以及流形是实的还是复的.下面两个结果表明，类似地，代数闭域上的簇的非奇异点处平展拓扑的局部环只依赖于簇的维数与域.

\begin{proposition}\label{prop:4.8}
若 $\varphi \colon Y \to X$ 在 $y$ 处平展，则 $\varphi$ 诱导的映射 $\mathcal{O}_{X, \bar{x}} \to \mathcal{O}_{Y, \bar{y}}$($x = \varphi(y)$)是同构.
\end{proposition}

\begin{proof}
把 $Y$ 换成 $y$ 的开邻域后，我们可以假定它在 $X$ 上平展(见命题~\ref{prop:2.14}).则 $y$ 的每个平展邻域也可以看作 $x$ 的平展邻域，而这样的邻域在 $x$ 的所有平展邻域中都是共尾的\footnote{预序集 $(A, \leq)$ 的子集 $B$ 是共尾的，如果对每个 $a \in A$，存在 $b \in B$ 使得 $a \leq b$.}.因此，两个正向极限相等.
\end{proof}

\begin{proposition}\label{prop:4.9}
设 $P$ 是簇 $X$ 的非奇异点，$d = \dim X$.则存在 $P$ 处平展的正则映射 $\varphi \colon U \to \mathbb{A}^{d}$.
\end{proposition}

\begin{proof}
选取在 $P$ 的邻域 $U$ 上定义的任意一组正则函数 $f_1, \ldots, f_d$，使得 $(df_1)_{P}, \ldots, (df_d)_{P}$ 是 $\operatorname{Tgt}_{P}(U)$ 的对偶空间的一组基(等价地，它们生成 $\mathcal{O}_{X, P}$ 中的极大理想).则
\[
\varphi \colon U \to \mathbb{A}^{d}, \quad \varphi(Q) = (f_1(Q), \ldots, f_d(Q)),
\]
在 $P$ 处平展，因为 $(d\varphi)_{P}$ 的对偶映射把 $(dX_i)$ 在 $\mathbb{A}^{d}$ 原点处的值送到 $(df_i)_{P}$.(更多细节见 AG \S 5.)
\end{proof}

注意该映射把 $P$ 送到 $\mathbb{A}^{d}$ 的原点.

\begin{proposition}\label{prop:4.10}
$\mathbb{A}^{d}$ 的原点处平展拓扑的局部环是
\[
k[[T_1, \ldots, T_d]] \cap k(T_1, \ldots, T_d)^{\mathrm{al}},
\]
即它由完备环 $k[[T_1, \ldots, T_d]]$ 中那些是 $k[T_1, \ldots, T_d]$ 中多项式(不一定是首一的)的根的元素组成.
\end{proposition}

我们将在下面讨论证明.

\subsection*{Hensel 环插曲}

\begin{proposition}\label{prop:4.11}
设 $A$ 是局部环，$k = A/\mathfrak{m}_{A}$.如下论断等价：
\begin{enumerate}
  \item $A$ 是 Hensel 的；
  \item 设 $f \in A[T]$；若 $\bar{f} = \bar{g_0}\bar{h_0}$，其中 $\bar{g_0}$ 首一且 $\bar{g_0}$ 与 $\bar{h_0}$ 互素，则 $f = gh$，其中 $g$ 首一，$\bar{g} = \bar{g_0}$, $\bar{h} = \bar{h_0}$；
  \item $A$ 满足上面的条件 ($\ast$)；
  \item 设 $B$ 是平展 $A$-代数；分解 $B/\mathfrak{m}_{A}B = k \times \bar{B_0}$ 提升为分解 $B = A \times B_0$.
\end{enumerate}
\end{proposition}

\begin{proof}
(d) $\Rightarrow$ (c).设 $f_1, \ldots, f_n$ 与 $a_0$ 如 ($\ast$) 中所述.设 $B = (A[T_1, \ldots, T_n]/(f_1, \ldots, f_n))_b$，其中 $b$ 已选取使得 $b(a_0) \neq 0$ 且 $\operatorname{Jac}(f_1, \ldots, f_n)$ 是 $B$ 中的单位.则 $A \to B$ 平展.由 $a_0$ 得到分解 $B/\mathfrak{m}_{A}B = k \times \bar{B_0}$，而 (d) 表明它提升为分解 $B = A \times B_0$.$n$ 元组 $(T_1, \ldots, T_n)$ 在投影映射 $B \to A$ 下的像就是 $a_0$ 到 $A$ 的提升.

(c) $\Rightarrow$ (b).引理~\ref{lem:4.7} 的证明可以修改以表明 ($\ast$) 蕴含 (b).

(b) $\Rightarrow$ (a).平凡.

(a) $\Rightarrow$ (d).我们可以假定 $A \to B$ 是标准平展同态(见 \S 2)，即 $B = (A[T]/(f(T)))_b$，其中 $f(T)$ 是首一多项式，使得 $f'(T)$ 在 $B$ 中可逆.分解 $B/\mathfrak{m}_{A}B = k \times \bar{B_0}$ 来自 $\bar{f}(T)$ 的单根 $a_0$，而 (a) 蕴含它提升为 $f(T)$ 在 $A$ 中的根.现在映射 $B \to A$，$T \mapsto a$，产生分解 $B = A \times B_0$.
\end{proof}

\begin{definition}\label{def:4.12}
设 $A$ 是局部环.如果每个其他局部同态 $A \to B$($B$ Hensel)都唯一分解为 $A \to A^{h} \to B$，则局部环的局部同态 $A \to A^{h}$($A^{h}$ Hensel)称为 $A$ 的 Hensel 化.
\end{definition}

显然，局部环的 Hensel 化是唯一的(如果存在).

\begin{proposition}\label{prop:4.13}
对局部环 $A$，
\[
A^{h} = \varinjlim_{(B, \mathfrak{q})} B,
\]
其中极限遍历所有这样的对 $(B, \mathfrak{q})$：$B$ 是平展 $A$-代数，$\mathfrak{q}$ 是素理想，使得 $\mathfrak{q} \cap A = \mathfrak{m}_{A}$，且诱导映射 $A/\mathfrak{m}_{A} \to B/\mathfrak{q}$ 是同构.
\end{proposition}

我们把证明留作习题.

\begin{corollary}\label{cor:4.14}
设 $X$ 是代数闭域 $k$ 上的簇.对每个 $x \in X$，$\mathcal{O}_{X, \bar{x}}$ 是 $\mathcal{O}_{X, x}$ 的 Hensel 化.
\end{corollary}

\begin{proposition}\label{prop:4.15}
设 $A$ 是局部环，$B$ 是所有满足
\[
A \subset H \subset \widehat{A}, \qquad \mathfrak{m}_{A} \subset \mathfrak{m}_{H} \subset \mathfrak{m}_{\widehat{A}}
\]
的局部 Hensel 环 $H$ 的交.则 $B$ 是 Hensel 的，且 $A \to B$ 是 $A$ 的 Hensel 化.
\end{proposition}

\begin{proof}
设 $f(T)$ 是 $B[T]$ 中的首一多项式，使得 $\bar{f} = \bar{g_0}\bar{h_0}$，其中 $\bar{g_0}$ 与 $\bar{h_0}$ 首一且互素.对每个 $H$，分解唯一提升为 $f = gh$，其中 $g, h \in H[T]$.由于唯一性，$g$ 与 $h$ 的系数在 $\bigcap H = B$ 中.因此，$B$ 是 Hensel 的.

设 $A \to A^{h}$ 是 $A$ 的 Hensel 化.因为 $B$ 是 Hensel 的，存在唯一的局部 $A$-同态 $A^{h} \to B$.该同态的像又是 Hensel 的，从而等于 $B$(因为按定义，$B$ 不含真的 Hensel 局部子环).
\end{proof}

M. Artin 有许多定理(和猜想)断言：已知在局部环的完备化上发生的事情，已经在 Hensel 化上发生.我们陈述其中之一.

\begin{theorem}\label{thm:4.16}
设 $k$ 是域，$f_1, \ldots, f_m \in k[X_1, \ldots, X_n, Y_1, \ldots, Y_r]$.设 $a_1, \ldots, a_r \in k[[X_1, \ldots, X_n]]$ 是形式幂级数，使得
\[
f_i(X_1, \ldots, X_n, a_1, \ldots, a_r) = 0, \qquad i = 1, \ldots, m.
\]
设 $N \in \mathbb{N}$.则存在 $k[X_1, \ldots, X_n]_{(X_1, \ldots, X_n)}$ 的 Hensel 化中的元素 $b_1, \ldots, b_r$，使得
\[
\begin{array}{l}
b_j \equiv a_j \pmod{(X_1, \ldots, X_n)^{N}}, \quad j = 1, \ldots, r; \\
f_i(X_1, \ldots, X_n, b_1, \ldots, b_r) = 0, \quad i = 1, \ldots, m.
\end{array}
\]
\end{theorem}

\begin{proof}
见 Artin 1973, p.48.
\end{proof}

\begin{corollary}\label{cor:4.17}
$A \overset{\mathrm{def}}{=} k[T_1, \ldots, T_d]_{(T_1, \ldots, T_d)}$ 的 Hensel 化是
\[
B \overset{\mathrm{def}}{=} k[[T_1, \ldots, T_d]] \cap k(T_1, \ldots, T_d)^{\mathrm{al}}.
\]
\end{corollary}

\begin{proof}
当然，按构造，$A^{h}$ 的每个元素都在 $A$ 上代数，从而 $A^{h} \subset B$.反过来，设 $\alpha \in k[[T_1, \ldots, T_d]]$ 是多项式 $f \in k[T_1, \ldots, T_d, X]$ 的根.对每个 $N \in \mathbb{N}$，存在 $a \in A^{h}$，使得 $f(a) = 0$ 且 $a \equiv \alpha \pmod{\mathfrak{m}^{N}_{A}}$.但 $f \in A[X]$ 只有有限多个根.因此，若对足够大的 $N$，$a \equiv \alpha \pmod{\mathfrak{m}^{N}_{A}}$，则 $a = \alpha$.
\end{proof}

局部环 $A$ 称为严格 Hensel 的，如果它是 Hensel 的且其剩余域可分代数闭，或者等价地，如果每个首一多项式 $f(T) \in A[T]$ 且 $\bar{f}(T)$ 可分，则它在 $A[T]$ 中分裂为次数 1 的因子.

\begin{definition}\label{def:4.18}
设 $A$ 是局部环.如果从 $A$ 到严格 Hensel 环 $H$ 的每个其他局部同态都延拓到 $A^{sh}$，而且一旦指定了剩余域上的映射 $A^{sh}/\mathfrak{m}^{sh} \to H/\mathfrak{m}_{H}$，延拓就唯一确定，则从 $A$ 到严格 Hensel 环 $A^{sh}$ 的局部同态 $A \to A^{sh}$ 是 $A$ 的严格 Hensel 化.
\end{definition}

例如，对域 $k$(看作局部环), Hensel 化就是 $k$ 本身，而 $k$ 的任何可分代数闭包都是 $k$ 的严格 Hensel 化.

\subsection*{概形}

概形 $X$ 的几何点是态射 $\bar{x} \colon \operatorname{Spec} \Omega \to X$，其中 $\Omega$ 是可分闭域.这样的点 $\bar{x}$ 的平展邻域是平展映射 $U \to X$ 连同位于 $\bar{x}$ 之上的几何点 $\bar{u} \colon \operatorname{Spec} \Omega \to U$.$\bar{x}$ 处平展拓扑的局部环是
\[
\mathcal{O}_{X, \bar{x}} \overset{\mathrm{def}}{=} \varinjlim_{(\bar{U}, \bar{u})} \Gamma(U, \mathcal{O}_U),
\]
其中极限遍历 $\bar{x}$ 的连通仿射平展邻域 $(U, \bar{u})$.

当 $X$ 是簇且 $x = \bar{x}$ 是 $X$ 的闭点时，这与前面的定义一致.

代数闭域上的簇的多数结果都\emph{照例}推广到概形.例如，$\mathcal{O}_{X, \bar{x}}$ 是 $\mathcal{O}_{X, x}$ 的严格 Hensel 化.

在本讲义余下部分，概形(或簇) $X$ 的几何点 $\bar{x}$ 处平展拓扑的局部环将称为 $\bar{x}$ 处的严格局部环，$\mathcal{O}_{X, \bar{x}}$.

代数概形的非闭点处的严格局部环在第 15 节(上同调维数)中使用.
